\documentclass[12pt]{report}

\usepackage{fancyhdr}
\setlength{\headheight}{15pt}
\usepackage{graphicx}   
\usepackage{subcaption}
\usepackage{extramarks}
\usepackage{amsmath}
\usepackage{amsthm}
\usepackage{amsfonts}
\usepackage[top=2cm, left=2cm, right=2cm, bottom=2.5cm]{geometry}
\usepackage{tikz}
\usepackage[plain]{algorithm}
\usepackage{algpseudocode}
\usepackage[colorlinks=true,linkcolor=black, citecolor=red]{hyperref}
\usepackage{url}
\usepackage{enumitem}
\usepackage{multirow}

\usepackage{listings}
\usepackage{xcolor} % required for coloring

\lstdefinestyle{pythonstyle}{
    language=Python,
    basicstyle=\ttfamily\small,
    keywordstyle=\color{blue}\bfseries,
    emphstyle=\color{red}\bfseries,
    stringstyle=\color{green!60!black},
    commentstyle=\color{gray}\itshape,
    numberstyle=\tiny\color{black!75},
    numbers=left,
    stepnumber=1,
    numbersep=8pt,
    showstringspaces=false,
    breaklines=true,
    frame=single,
    rulecolor=\color{black!75},
    backgroundcolor=\color{white},
    tabsize=4,
    captionpos=b,
}
\lstset{style=pythonstyle}

\usepackage{nomencl}
\renewcommand{\nomname}{List of Symbols}
\makenomenclature

\usepackage{scalerel}

\setlength{\parindent}{0mm}
\setlength{\parskip}{2mm}

\usetikzlibrary{automata,positioning}

\pagestyle{fancy}

\lhead{\reportTitle}
\rhead{\authorName \textbf{ - \the\year{}}}


\cfoot{\thepage}

% \renewcommand\headrulewidth{0.4pt}
% \renewcommand\footrulewidth{0.4pt}

\newcommand{\reportTitle}{Finite Elements for Linear and Non Linear Mechanics}
\newcommand{\authorName}{\textbf{Gabriel Jezler}}

%
% Title Page
%

\title{
    \vspace{5cm}
    \textmd{\textbf{\reportTitle}}\\
    \vspace{8cm}
}

\author{\authorName}
\date{\today}

% \renewcommand{\part}[1]{\textbf{\large Part \Alph{partCounter}}\stepcounter{partCounter}\\}

%
% Various Helper Commands
%
% For derivatives
\newcommand{\deriv}[1]{\frac{\mathrm{d}}{\mathrm{d}x} (#1)}

% For partial derivatives
\newcommand{\psderiv}[2]{\frac{\partial}{\partial #2} (#1)}
\newcommand{\pderiv}[2]{\frac{\partial #1}{\partial #2} }

% Integral dx
\newcommand{\dt}{\, \mathrm{dt}}
\newcommand{\dV}{\, \mathrm{dV}}
\newcommand{\dS}{\, \mathrm{dS}}
\newcommand{\dv}{\, \mathrm{dv}}
\newcommand{\ds}{\, \mathrm{ds}}

% Tensor
\newcommand{\tensorTwo}[1]{\bar{\bar{#1}}}
\newcommand{\tensorFour}[1]{\bar{\bar{\bar{\bar{#1}}}}}

\newcommand{\tensorTwoGradO}[1]{\tensorTwo{\nabla}_\mathrm{0} #1}
\newcommand{\vectorGradO}[1]{\vec{\nabla}_\mathrm{0} #1}
\newcommand{\tensorTwoGrad}[1]{\tensorTwo{\nabla} #1}
\newcommand{\vectorGrad}[1]{\vec{\nabla} #1}

% Kinematics
\newcommand{\x}[0]{\vec{x}}
\newcommand{\X}[0]{\vec{X}}

\newcommand{\F}[0]{\tensorTwo{\mathrm{F}}}
\newcommand{\C}[0]{\tensorTwo{\mathrm{C}}}
\newcommand{\E}[0]{\tensorTwo{\mathrm{E}}}

\newcommand{\J}[0]{\mathrm{J}}

% Stress
\newcommand{\cauchy}[0]{\tensorTwo{\sigma}}
\newcommand{\kirchhoff}[0]{\tensorTwo{\scalerel*{\tau}{X}}}
\newcommand{\pkOne}[0]{\tensorTwo{\mathrm{P}}}
\newcommand{\pkTwo}[0]{\tensorTwo{\mathrm{S}}}

% Potential
\newcommand{\potential}[0]{\Psi}

% List of Symbols
\nomenclature{$\X$}{Position vector at the initial configuration}
\nomenclature{$\vec{\phi}$}{Transformation function}
\nomenclature{$\x$}{Position vector at the deformed configuration}
\nomenclature{$\F$}{Transformation gradient tensor}
\nomenclature{$\J$}{Jacobian of the transformation}
\nomenclature{$\C$}{Cauchy-Green right tensor}
\nomenclature{$\E$}{Green-Lagrange tensor}
\nomenclature{$\cauchy$}{Cauchy stress tensor}
\nomenclature{$\kirchhoff$}{Kirchhoff stress tensor}
\nomenclature{$\pkOne$}{First Piola-Kirchhoff stress tensor}
\nomenclature{$\pkTwo$}{Second Piola-Kirchhoff stress tensor}
\nomenclature{$\Psi$}{Strain energy density}
\nomenclature{$\tensorTwo{\varepsilon}$}{Strain tensor}
\nomenclature{$\delta \mathcal{W}$}{Virtual work of the system}
\nomenclature{$\vec{n}$}{Normal vector}
\nomenclature{$v$}{Volume of the deformed configuration}
\nomenclature{$V$}{Volume of the initial configuration}
\nomenclature{$s$}{Surface of the deformed configuration}
\nomenclature{$S$}{Surface of the initial configuration}
\nomenclature{$N$}{Shape function}

\begin{document}

\maketitle

\pagebreak

\printnomenclature

\pagebreak

\tableofcontents

\pagebreak

\chapter{Introduction}

This document aims to provide a theoretical explanation for the finite element method (FEM) for solid mechanics problems. The document begins by presenting all the necessary theoretical background, including the definition of many kinematics and stress related variables. From those, the available material constitutive models are presented, which are used to relate the kinematics and stress variables. Then, equation to be solved is derived by using the principle of virtual work. It is then discretized by using the finite element method. Finally, both solution methods are explained. Addionally, the document will present an overview of the code structure and all of its components.

This project is based on a academic projet I made in 2025 at Centrale Nantes. Based on that, the approach has been adapted to be used in a more general way, and the code has been rewritten for better comprehension, modularity and optimize performance - obtained mainly through the use of numpy array instead of loops. Addionally, the book \textit{"Nonlinear Continuum Mechanics for Finite Element Analysis"} by Javier Bonet and Richard D. Wood was used for theoretical reference.


\chapter{Kinematics and Stresses}

In this section, the kinematics and the different stress tensors used in this paper are defined. The relation between them is also given. Finally, the definition of the strain energy density is given as well as its relation with the different stress tensors. These definitions are used throughout the paper and are important for the understanding of the following sections.

\section{Kinematics}

The kinematics of a body is defined by the transformation between the initial configuration and the current configuration. This transformation can be defined as stated by Equation \ref{eq:transformation}. Thoughout this paper, the initial configuration will be denoted by capital letters while the current configuration will be denoted by lowercase letters.

\begin{equation} \label{eq:transformation}
    \x = \vec{\phi} (\X, \, t)
\end{equation}

Based on this transformation, it is possible to define the transformation gradient (Equation \ref{eq:F}) and the Jacobian of the transformation (Equation \ref{eq:J}).

\begin{equation} \label{eq:F}
    \F = \tensorTwoGradO{\x} = \pderiv{\x}{\X}
\end{equation}

\begin{equation} \label{eq:J}
    \J = \det{\F}
\end{equation}

Based on the tensor $\F$, it is possible to define multiple tensors to represent the defromation. In this codes, two of then are used: the right Cauchy-Green deformation tensor $\C$ and the Green-Lagrange strain tensor $\E$ as given by Equations \ref{eq:C} and \ref{eq:E} respectively.

\begin{equation} \label{eq:C}
    \C = \F^T \F
\end{equation}

\begin{equation} \label{eq:E}
    \E = \frac{1}{2} \left(\C - \tensorTwo{\mathrm{I}}\right) = \frac{1}{2} \left(\F^T \F - \tensorTwo{\mathrm{I}}\right)
\end{equation}

\section{Stresses}

Four different stress tensors are used throughout this paper and are listed below:

\begin{itemize}
    \item $\cauchy$: Cauchy stress tensor
    \item $\kirchhoff$: kirchhoff stress tensor
    \item $\pkOne$: First Piola-Kirchhoff stress tensor
    \item $\pkTwo$: Second Piola-kirchhoff stress tensor
\end{itemize}

They are all stress' representations but under different configurations. They can all be related under the following equation:

\begin{equation}
    \cauchy = \J^{-1} \cdot \kirchhoff = \J^{-1} \cdot \pkOne \F^T = \J^{-1} \cdot \F \pkTwo \F^T
\end{equation}

\section{Strain Energy Density}

The strain energy density is defined as:

\begin{equation}
    \potential (\F, \, \X) = \int\pkOne :\dot{\F} \dt
\end{equation}

From this definition, a new definition for $\pkOne$ can be stated:

\begin{equation}
    \pkOne = \pderiv{\potential}{\F}
\end{equation}

This can also be written as a function of the tensor $\C$ as given below:

\begin{equation}
    \potential (\F, \, \X) = \potential (\C, \, \X) = \int \frac{1}{2}\pkTwo : \dot{\C} \dt 
\end{equation}

This means the $\pkTwo$ can be written as:

\begin{equation}
    \pkTwo = 2 \pderiv{\Psi}{\C} = \pderiv{\Psi}{\E}
\end{equation}


\chapter{Material Behavior}

\section{Linear Materials}

\subsection{Linear Elastic Isotropic} \label{mat:linear_elastic_iso}

This model uses the hypothesis of small displacements under a elastic regime. This is mainly done by assuming that:

\begin{equation*}
    \tensorTwoGrad{\vec{u}} \approx \tensorTwoGradO{\vec{u}} \qquad \text{and} \qquad \lVert  \tensorTwoGrad{\vec{u}} \lVert \approx 0
\end{equation*}

This leads to the following conclusions:

\begin{equation*}
    \left\{
    \begin{array}{cc}
        \F &= \tensorTwo{\mathrm{I}} + \tensorTwoGradO{\vec{u}} \approx \tensorTwo{\mathrm{I}} \\
        \J &= \det{\F} \approx 1 \\
    \end{array} \right.
    \qquad \Longrightarrow \qquad \cauchy \approx \kirchhoff \approx \pkOne \approx \pkTwo
\end{equation*}

Under this hypothesis, the strain ($\tensorTwo{\varepsilon}$) is used to model the material. It is defined as:

\begin{equation} \label{eq:strain}
    \tensorTwo{\varepsilon} = \frac{1}{2} \left(\tensorTwoGrad{\vec{u}} + \tensorTwoGrad{\vec{u}}^T \right)
\end{equation}

Using this, both the strain energy density and Cauchy stress can be defined as:

\begin{equation} \label{eq:psi_linear_elastic_iso}
    \potential = \frac{1}{2} \lambda \, \left( tr(\tensorTwo{\varepsilon}) \right)^2 + \mu \,\tensorTwo{\varepsilon} : \tensorTwo{\varepsilon}
\end{equation}

\begin{equation} \label{eq:sigma_linear_elastic_iso}
    \cauchy = \lambda \, tr(\tensorTwo{\varepsilon}) \tensorTwo{\mathrm{I}} + 2\mu \,\tensorTwo{\varepsilon}
\end{equation}

The Equation \ref{eq:sigma_linear_elastic_iso} is also known as the generalized Hooke's law for isotropic materials. This is written as a function of the elaticity tensor, which is given by:

\begin{equation} \label{eq:sigma_linear_elastic_iso_C}
    \cauchy = \tensorFour{C} : \tensorTwo{\varepsilon} \qquad \text{with} \qquad
    C_{ijkl} = \lambda \, \delta_{ij}\delta_{kl} + \mu \left( \delta_{ik}\delta_{jl} + \delta_{il}\delta_{jk} \right)
\end{equation}

This model is characterized by two parameters ($\lambda$ and $\mu$) which are called the Lamé coefficients. They can be written as a function of the Young's modulus ($E$) and Poisson coefficient ($\nu$) as follows:

\begin{equation} \label{eq:lame1}
    \lambda = \frac{E \nu}{(1 + \nu)(1 - 2\nu)}
\end{equation}
\begin{equation} \label{eq:lame2}
    \mu = \frac{E}{2(1 + \nu)}
\end{equation}

\section{Non Linear Materials}

\subsection{Saint Venant-Kirchhoff} \label{mat:svt}

This model can be is often used for moderate deformations under a elastic regime. Unlike the \nameref{mat:linear_elastic_iso} model, the present framework does not uses the small displacement hypothesis and therefore is no longer a linear model.

Similarly, the model uses the strain energy density to correlate the stress and kinematics of the material. Within this context, the strain energy and the second Piola-Kirchhoff stress tensor are give below:

\begin{equation} \label{eq:psi_svt}
    \potential = \frac{1}{2} \lambda \, \left( tr(\E) \right)^2 + \mu \, \E : \E
\end{equation}

\begin{equation} \label{eq:pk2_svt}
    \pkTwo = \lambda \, tr(\E) \tensorTwo{\mathrm{I}} + 2\mu \,\E
\end{equation}

\subsection{Compressible Neo-Hookean} \label{mat:nh}

Unlike the \nameref{mat:svt} model, this model can be employed for large deformations under a elastic regime. It considers a compressible hyper-elastic material that follows the following equations:

\begin{equation} \label{eq:psi_nh}
    \potential = \frac{1}{2} \mu \, \left( tr(\C) - 3 \right) + \mu \, \ln{\J} + \frac{1}{2} \lambda \, \left( \ln{\J} \right)^2
\end{equation}

\begin{equation} \label{eq:pk2_nh}
    \pkTwo = \mu \, \left(\tensorTwo{\mathrm{I}} - \C^{-1} \right) - \lambda \, \ln{\J} \, \C^{-1}
\end{equation}

\chapter{Finite Element Formulation}

\section{Equilibrium Equation}

The equation to be solved is given below:

\begin{equation}
    \vec{\mathrm{div}} \cauchy + \vec{f} = \vec{0}
\end{equation}

However, the equation above is in its strong formulation. For the finite elements methods, it is far more useful to convert it to its weak form. This allows for the application of boundary conditions directly in the equation. It can be done by using the virtual work theorem:

\begin{equation*}
    \vec{\mathrm{div}} \cauchy + \vec{f} = \vec{0} \qquad \Longrightarrow \qquad 
    \delta \mathcal{W} = \int_v \left( \vec{\mathrm{div}} \cauchy + \vec{f} \right) \cdot \vec{\delta u} \dv = 0 \quad \forall \, \vec{\delta u} \Longrightarrow
\end{equation*}
\begin{equation*}
    \Longrightarrow \qquad \delta \mathcal{W} = \int_v \vec{\mathrm{div}} \cauchy \cdot \vec{\delta u} \dv + \int_v \vec{f} \cdot \vec{\delta u} \dv = 0 \quad \forall \, \vec{\delta u}
\end{equation*}

This can be further developed by applying the following rule (Equation \ref{eq:div_prod_rule}) along side with the divergence theorem (Equation \ref{eq:divergence_th}).

\begin{equation} \label{eq:div_prod_rule}
    \vec{\mathrm{div}} \tensorTwo{A} \cdot \vec{b} = \vec{\mathrm{div}} \left( \tensorTwo{A} \cdot \vec{b} \right) - \tensorTwo{A} : \vec{b}
\end{equation}

\begin{equation} \label{eq:divergence_th}
    \int_v \mathrm{div} \left(\vec{a} \right) \dv = \int_s \vec{a} \cdot \vec{n} \ds
\end{equation}

By applying both relation on the left hand term of the equation, it is possible to obtain:

\begin{equation} \label{eq:wf_cauchy}
    \delta \mathcal{W} = \int_v \cauchy : \tensorTwoGrad{\vec{\delta u}} \dv - \left (\int_v \vec{f} \cdot \vec{\delta u} \dv + \int_s \vec{t} \cdot \vec{\delta u} \ds \right) = 0 \quad \forall \, \vec{\delta u}
\end{equation}

Where:

\begin{equation*}
    \vec{t} = \cauchy \cdot \vec{n}
\end{equation*}

The next step is to rewrite Equation \ref{eq:wf_cauchy} with respect to the initial configuration. This can be done by using the Jacobian of the transformation defined by Equation \ref{eq:J}. Using this, it is possible to rewrite Equation \ref{eq:wf_cauchy} as follows:

\begin{equation*}
    \int_V \cauchy : \tensorTwoGrad{\vec{\delta u}} \, J \dV - \left (\int_V \vec{f} \cdot \vec{\delta u} \, J \dV + \int_S \vec{t} \cdot \vec{\delta u} \left( \pderiv{s}{S} \right)\dS \right) = 0 \quad \forall \, \vec{\delta u} \qquad \Longrightarrow
\end{equation*}
\begin{equation*}
    \Longrightarrow \qquad \int_V \underbrace{\kirchhoff}_{\J \cauchy} : \tensorTwoGrad{\vec{\delta u}} \dV - \left (\int_V \underbrace{\vec{f_0}}_{J \vec{f}} \cdot \vec{\delta u} \dV + \int_S \underbrace{\vec{t_0}}_{\left( \pderiv{s}{S} \right) \vec{t}} \cdot \vec{\delta u} \dS \right) = 0 \quad \forall \, \vec{\delta u} 
\end{equation*}

Furthermore, the gradient can be written as:

\begin{equation*}
    \tensorTwoGrad{\vec{\delta u}} = \pderiv{\vec{\delta u}}{\x} = \pderiv{\vec{\delta u}}{\X} \cdot \pderiv{\X}{\x} = \tensorTwoGradO{\vec{\delta u}} \cdot \F^{-1}
\end{equation*}

Finally, using that:

\begin{equation*}
    \tensorTwo{A} : \left( \tensorTwo{B} \tensorTwo{C} \right) = \left( \tensorTwo{A} \tensorTwo{C}^T \right) : \tensorTwo{B}
\end{equation*}

It is possible to obtain the final equation:

\begin{equation} \label{eq:wf_pk1}
    \delta \mathcal{W} = \int_V \underbrace{\pkOne}_{\kirchhoff \F^{-T}} : \tensorTwoGradO{\vec{\delta u}} \dV - \left (\int_V \vec{f_0} \cdot \vec{\delta u} \dV + \int_S \vec{t_0} \cdot \vec{\delta u} \dS \right) = 0 \quad \forall \, \vec{\delta u} 
\end{equation}

\section{Discretization}

The discretization is a important step for the finite volume method. In this part, we use interpolation based on nodal values for the value of displacement inside an element. This is done by using shape functions ($N$) for each element node $a$ ($N_a$).

\begin{equation}
    \vec{u} ( \X ) = \sum_{a \, \in \, e} N_a (\vec{\xi}) \cdot \vec{U}_a
\end{equation}
\begin{equation}
    \tensorTwoGradO{\vec{u}}( \X ) = \sum_{a \, \in \, e} \tensorTwoGradO{N_a}(\vec{\xi}) \otimes \vec{U}_a
\end{equation}

In those, the shape function is given as a function of the position vector $\vec{\xi}$. This is due to the fact that these functions are defined for a reference coordinate space. For a line element this would be for $\xi \in [-1, 1]$, for example. It is then possible to use the element's transformation (Equation \ref{eq:elem_transf}) to relate the reference system ($\vec{\xi}$) to the element's global coordinates ($\X$).

\begin{equation} \label{eq:elem_transf}
    \X = \psi (\vec{\xi}) \qquad \Longrightarrow \qquad \vec{\xi} = \psi^{-1} (\X)
\end{equation}

In addition, the Galerkin method is used, which means that the virtual field is given by:

\begin{equation}
    \vec{\delta u} ( \X ) = \sum_{a \, \in \, e} N_a (\vec{\xi}) \cdot \vec{\delta U}_a
\end{equation}
\begin{equation}
    \tensorTwoGradO{\vec{\delta u}}( \X ) = \sum_{a \, \in \, e} \tensorTwoGradO{N_a}(\vec{\xi}) \otimes \vec{\delta U}_a
\end{equation}

\section{Elementary description}

The total virtual work of the system can be given as the sum of the virtual work for each element as stated below. This means that the final equation will be the result of the sum of the contribution of each element. For that reason, it is possible to work with the element-wise contribution and then add them up in a process called assembly.

\begin{equation}
    \delta \mathcal{W} = \sum_e \delta \mathcal{W}^{(e)} = 0 \quad \forall \, \vec{\delta u}
\end{equation}

Therefore, by analyzing the element contribution fro the virtual work, the following relation is obtained

\begin{equation*}
    \delta \mathcal{W}^{(e)} = \int_{V^{(e)}}\pkOne^{(e)} : \tensorTwoGradO{\vec{\delta u}} \dV - \left (\int_{V^{(e)}} \vec{f_0^{(e)}} \cdot \vec{\delta u} \dV + \int_{S^{(e)}} \vec{t_0^{(e)}} \cdot \vec{\delta u} \dS \right)
\end{equation*}

 By using the discretized functions in the equation above, it then becomes:

\begin{equation} \label{eq:dW_disc}
    \delta \mathcal{W}^{(e)} = \sum_{a \, \in \, e} \left[ \underbrace{\int_{V^{(e)}} \pkOne^{(e)} \vectorGradO{N_a} \dV}_{\vec{\delta \mathcal{W}}_{int, a}^{(e)}} - \underbrace{\left (\int_{V^{(e)}} \vec{f_0^{(e)}} \cdot N_a  \dV + \int_{S^{(e)}} \vec{t_0^{(e)}} \cdot N_a \dS \right)}_{\vec{\delta \mathcal{W}}_{ext, a}^{(e)}} \right] \cdot \vec{\delta U}_a
 \end{equation}

\section{Numerical Integration}

To solve for the integrals in the equations above, a numerical integration strategy is used. This is done with the Gauss quadrature framework. 

Before using this process, one additional development needs to be done to Equation \ref{eq:dW_disc}. As mentioned before, the shape functions are written in the space of the reference element. Therefore, it is useful to also integrate in this coordinate system. To make this change in integration variable, a volume and surface Jacobian for the $V \longrightarrow \hat V$ and $S \longrightarrow \hat S$ transformations are used and are represented by $\J_V$ and $\J_S$, respectively. 

\begin{align*}
    \delta \mathcal{W}_{int,a}^{(e)} &= \int_{V^{(e)}} \pkOne^{(e)} \vectorGradO{N_a} \dV = \\
    &= \int_{\hat V^{(e)}} \pkOne^{(e)} (\psi^{(e)}(\vec{\xi})) \vectorGradO{N_a} (\vec{\xi}) \cdot \J_V^{(e)} \mathrm{d\hat V} \\
    \delta \mathcal{W}_{ext,a}^{(e)} &= \int_{V^{(e)}} \vec{f_0^{(e)}} \cdot N_a  \dV + \int_{S^{(e)}} \vec{t_0^{(e)}} \cdot N_a \dS = \\
    &= \int_{\hat{V}^{(e)}} \vec{f_0^{(e)}} (\psi^{(e)}(\vec{\xi})) \cdot N_a (\vec{\xi}) \cdot \J_V^{(e)} \mathrm{d\hat V} + \int_{\hat{S}^{(e)}} \vec{t_0^{(e)}} (\psi^{(e)}(\vec{\xi})) \cdot N_a (\vec{\xi}) \cdot \J_S^{(e)} \mathrm{d\hat S}
\end{align*}

Now, by using the quadrature's approximation, it is possible to write:

\begin{equation} \label{eq:integral_Wint}
    \delta \mathcal{W}_{int,a}^{(e)} \approx \sum_g \omega_g \cdot  \pkOne^{(e)} (\psi^{(e)}(\vec{\xi_g})) \vectorGradO{N_a} (\vec{\xi_g}) \cdot \J_V^{(e)}
\end{equation}
\begin{equation} \label{eq:integral_Wext}
    \delta \mathcal{W}_{ext,a}^{(e)} \approx \sum_g \omega_g \cdot \vec{f_0^{(e)}} (\psi^{(e)}(\vec{\xi_g})) \cdot N_a (\vec{\xi_g}) \cdot \J_V^{(e)} + \sum_{g_s} \vec{t_0^{(e)}} (\psi^{(e)}(\vec{\xi_{g_s}})) \cdot N_a (\vec{\xi_{g_s}}) \cdot \J_S^{(e)}
\end{equation}


\chapter{Linear Solution}

This chapter is dedicated to the solution of the linear system of equations obtained from the finite element formulation. This approach is used to the linear materials presented in Section \ref{sec:lin_mat}. The final system of equations has the form given by Equation \ref{eq:linear-form} \footnote{Note that the linear system is expressed with tensor and vector notation even if, technically, they are (N$\times$N) and (N$\times$1) matrices, respectively.}.

\begin{equation} \label{eq:linear-form}
    \tensorTwo{K} \cdot \vec{U} = \vec{F}
\end{equation}

In this case, $\tensorTwo{K}$ is the stiffness matrix, $\vec{U}$ nodal displacement vector - unknown - and $\vec{F}$ is the vector of external forces. In order to obtain this system, the Equation \ref{eq:assembled_eq} will serve as a starting point. The internal virtual work term ($\delta \vec{\mathcal{W}}_{int, a}^{(e)}$) will be used to generate the $\tensorTwo{K} \cdot \vec{U}$ term, while the external virtual work term ($\delta \vec{\mathcal{W}}_{ext, a}^{(e)}$) will be used to generate the $\vec{F}$ term.

To obtain the stiffness matrix, we need to separate the nodal displacements from the rest of the terms in order to obtain the $\tensorTwo{K}$ and $\vec{U}$ vector. In order to do this, the linearity of the material along with the hypothesis of small displacements (Equation \ref{eq:small_disp}) are used. We will, first, developthe stress - strain relation:

\begin{equation*}
    \cauchy = \tensorFour{\mathcal{C}} : \tensorTwo{\varepsilon} \Longrightarrow \cauchy_{ij} = \mathcal{C}_{ijkl} \cdot \varepsilon_{kl}
\end{equation*}

Additionally, from the definition of strain (Equation \ref{eq:strain}) and the finite element discretization (Equation \ref{eq:disc_grad0_u}), the following relation is obtained:

\begin{equation*}
    \varepsilon_{kl} = \frac{1}{2} \left(\tensorTwoGrad{\vec{u}}_{kl} + \tensorTwoGrad{\vec{u}}_{lk} \right) = \frac{1}{2} \left(\vectorGradO{N_b} \left(\vec{\xi} \right) \otimes \vec{U}_b + \vec{U}_b \otimes \vectorGradO{N_b} \left(\vec{\xi} \right)  \right) \Longrightarrow
\end{equation*}
\begin{equation}
    \varepsilon_{kl} = \frac{1}{2} \left(\nabla_\mathrm{0} N_{bk} \left(\vec{\xi} \right) \cdot \delta_{lm} + \nabla_\mathrm{0} N_{bl} \left(\vec{\xi} \right) \cdot \delta_{km}  \right) \cdot U_{bm} = B_{bmkl} \left(\vec{\xi} \right) \cdot U_{bm}
\end{equation}

So, it is possible to use Equation \ref{eq:integral_Wint} to obtaint the stiffness matrix:

\begin{equation*}
    \delta {\mathcal{W}}_{int,ai}^{(e)} \approx \left[ \sum_g \omega_g \cdot   \mathcal{C}_{ijkl} \cdot B_{bmkl} \left(\vec{\xi}_g \right) \cdot \nabla_\mathrm{0} N_{aj} \left(\vec{\xi}_g \right) \cdot \J_V^{(e)} \right] \cdot U_{bm} \Longrightarrow
\end{equation*}
\begin{equation}
    \delta {\mathcal{W}}_{int,ai}^{(e)} \approx K_{aibm}^{(e)} \cdot U_{bm} \qquad \text{where} \qquad K_{aibm}^{(e)} = \sum_g \omega_g \cdot \mathcal{C}_{ijkl} \cdot B_{bmkl} \left(\vec{\xi}_g \right) \cdot \nabla_\mathrm{0} N_{aj} \left(\vec{\xi}_g \right) \cdot \J_V^{(e)}
\end{equation}

Additionally, the external force vector is obtained by summing the contribution of each element to the external virtual work, which is done by the expression in Equation \ref{eq:integral_Wext}. Once the assembly process is done, the following system will be solved \footnote{It is important to note that the indices $a, m and b, i$ represents the number of nodes and dimensions of the problem, respectively. To solve it numerically, the solver will need to combine them into a single index: $n_{nodes} \cdot dim$.}:

\begin{equation}
    K_{aibm} \cdot U_{bm} = F_{ai} 
\end{equation}


\chapter{Non Linear Solution}

\section{Newton-Raphson Method}

The Newton-Raphson method is an iterative method used to solve non linear equations. For the case of a 1D equation, the formulation is given by:

\begin{equation*}
    f(x) = 0
\end{equation*}

The ideia is to start with an initial guess for the solution and then iteratively improve it until convergence is achieved. The method is based on the Taylor series expansion of the function around the current guess. The next guess is then given by:

\begin{equation} \label{eq:newton-raphson-1d}
    x^{k+1} = x^k + \Delta x^k \qquad \text{where} \qquad \Delta x^k = -\frac{f(x^k)}{f'(x^k)}
\end{equation}

As we expand this formulation, it is possible to solve for a general $f$ that can be a scalar, vector, tensor or any type of function with a vector $\vec{x}$ as input:

\begin{equation*}
    f \left(\vec{x} \right) = 0
\end{equation*}

For this case, the definition of a directional derivative is needed. This can be defined as:

\begin{equation} \label{eq:directional-der}
    \mathrm{D} f \left( \vec{x} \right) \, \left[\vec{u} \right] = \frac{d}{d \epsilon} \Bigr \rvert_{\epsilon = 0} f \left(\vec{x} + \epsilon \vec{u} \right)
\end{equation}

Therefore, the next guess for the solution is given by:

\begin{equation} \label{eq:newton-raphson}
    x^{k+1} = x^k + \Delta x^k \qquad \text{where} \qquad \mathrm{D} f \left( \vec{x}^k \right) \, \left[\Delta \vec{x}^k \right] = -f \left(\vec{x}^k \right)
\end{equation}

\section{Application to Finite Elements}

In order to apply the Newton-Raphson method to the finite element method, the equation to be solved is stated in \ref{eq:wf_pk2}. This means that the function to be solved is the virtual work of the system, and hence, the firts step is to compute the directional derivative of the virtual work with respect to a displacement increment $\Delta \vec{u}$:

\begin{equation} \label{eq:dir-der-W}
    \begin{split}
        &\mathrm{D} \delta \mathcal{W} \left( \vec{\phi}, \, \delta \vec{u} \right) \, \left[\Delta \vec{u} \right] = \\\\
        &= \underbrace{\int_V \delta \Edot : \mathrm{D} \pkTwo [\Delta \vec{u}] \dV + \int_V \pkTwo : \mathrm{D} \delta \Edot [\Delta \vec{u}] \dV}_{\mathrm{D} \delta \mathcal{W}^{int} \, \left[\Delta \vec{u} \right]} - \underbrace{\left(\int_V \mathrm{D} \vec{f_0} \, [\Delta \vec{u}] \cdot \delta \vec{v} \dV + \int_S \mathrm{D} \vec{t_0} \, [\Delta \vec{u}] \cdot \delta \vec{v} \dS \right)}_{\mathrm{D} \delta \mathcal{W}^{ext} \, \left[\Delta \vec{u} \right]} 
    \end{split}
\end{equation}

To solve for the next step, the virtual work of the system is composed of the sum of the internal virtual work and the external virtual work. Then, the directional derivative of each of them is computed separately. However, the directional derivative of the virtual work has the following form \footnote{Note that in this section the vectors $\delta \vec{v}_a$ and $\Delta \vec{U}_b$ are writen with eachs of their components : $\delta \vec{v}_a = \delta v_{aj} \cdot \hat{e}_j$ and $\Delta \vec{U}_b = \Delta U_{bl} \cdot \hat{e}_l$.}:

\begin{equation*}
    D \delta \mathcal{W} \, \left[\Delta \vec{u} \right] = \delta V_{aj} \cdot K_{ajbl} \cdot \Delta U_{bl}
\end{equation*}

The goal with this section is to find the tangent tensor $K_{ajbl}$ for each combination of nodes $a$ and $b$ for both internal and external virtual work contributions. Once this is done, the next step is to solve for the increment of displacement $\Delta \vec{U}$ by solving the following system of equations:

\begin{equation*}
    D \delta \mathcal{W} \, \left[\Delta \vec{u} \right] = - \delta \mathcal{W} \qquad \Longrightarrow \qquad \delta V_{aj} \cdot K_{ajbl} \cdot \Delta U_{bl} = \left[-R_{aj} \right] \cdot \delta V_{aj} \qquad \forall \, \delta V_{aj} \qquad \Longrightarrow
\end{equation*}
\begin{equation}
    \qquad \underbrace{\left[K_{ajbl} \cdot \Delta U_{bl} + R_{aj} \right]}_{0} \cdot \delta V_{aj} = 0 \qquad \forall \, \delta V_{aj}
\end{equation}

\subsection{Internal virtual work}

The goal of this section is to compute the directional derivative of the internal virtual work. To achieve this, it is first necesary to compute the directional derivative of the second Piola-Kirchhoff stress tensor $\pkTwo$ and the virtual deformation rate tensor $\delta \Edot$. For the case of $\delta \Edot$, it is possible to use the definition of the deformation rate tensor and the directional derivative to obtain:

\begin{equation} \label{eq:der-dir-Edot}
    D \delta \Edot \, [\Delta \vec{u}] = D \left( \frac{1}{2} \left( \delta \Fdot^T \cdot \F + \F^{T} \cdot \delta \Fdot  \right) \right) [\Delta \vec{u}] = \frac{1}{2} \left( \delta \Fdot^T \cdot \tensorTwoGradO{\left( \Delta \vec{u} \right)} + \tensorTwoGradO{\left( \Delta \vec{u} \right)}^{T} \cdot \delta \Fdot  \right)
\end{equation}


Additionally, the same can be done for the directional derivative of $\pkTwo$:

\begin{equation} \label{eq:der-dir-pk2}
    \mathrm{D} \pkTwo \, [\Delta \vec{u}] = \Celastic : \mathrm{D} \E \, [\Delta \vec{u}] = \Celastic : \frac{1}{2} \left( \F^T \cdot \tensorTwoGradO{\left( \Delta \vec{u} \right)} + \tensorTwoGradO{\left( \Delta \vec{u} \right)}^{T} \cdot \F  \right)
\end{equation}

Then, using that into Equation \ref{eq:dir-der-W} and the symmetries of $\pkTwo$ and $\Celastic$, the following expression is obtained for the directional derivative of the internal virtual work:

\begin{equation}
    D \delta \mathcal{W}^{int} \, \left[\Delta \vec{u} \right] = \int_V \left( \tensorTwoGradO{\delta \vec{v}} \cdot \F \right) : \Celastic : \left( \F \cdot \tensorTwoGradO{\left( \Delta \vec{u} \right)} \right) \dV + \int_V \pkTwo : \left( \tensorTwoGradO{\left( \Delta \vec{u} \right)}^{T} \cdot \tensorTwoGradO{\delta \vec{v}} \right) \dV
\end{equation}

With this expression it is possible to isolate the tangent tensor $\tensorFour{K}^{int}$. This is done by using the definition of the finite element discretization of the displacement field. The final expression for the tangent tensor $\tensorFour{K}^{int}$ is given by:

\begin{align*}
    D \delta \mathcal{W}^{int} \, \left[\Delta \vec{u} \right] &= \int_V \left( \tensorTwoGradO{\delta \vec{v}} \cdot \F \right) : \Celastic : \left( \F \cdot \tensorTwoGradO{\left( \Delta \vec{u} \right)} \right) \dV + \int_V \pkTwo : \left( \tensorTwoGradO{\left( \Delta \vec{u} \right)}^{T} \cdot \tensorTwoGradO{\delta \vec{v}} \right) \dV \\
    &= \int_V \left[ \left(\vectorGradO{N_a} \otimes \delta \vec{V}_{a} \right) \cdot \F \right] : \Celastic : \left[ \F \cdot \left(\vectorGradO{N_b} \otimes \Delta \vec{U}_{b} \right) \right] \dV + \\
    & \qquad + \int_V \pkTwo : \left[ \left( \Delta \vec{U}_{b} \otimes \vectorGradO{N_b} \right) \cdot \left(\vectorGradO{N_a} \otimes \delta \vec{V}_{a} \right) \right] \dV \\
\end{align*}
\begin{equation}
    \begin{split}
        D \delta \mathcal{W}^{int} \, \left[\Delta \vec{u} \right] &= \delta V_{aj} \cdot \Bigg[ \underbrace{\int_V (\nabla_0 N)_{ai} \cdot F_{jk} \cdot \mathcal{C}_{ikml} \cdot F_{mn} \cdot (\nabla_0 N)_{bn} \dV}_{K_{ajbl}^{int, 1}} + \\
        & \qquad + \underbrace{\int_V (\nabla_0 N)_{ak}\cdot \mathrm{S}_{jl} \cdot (\nabla_0 N)_{bk} \dV}_{K_{ajbl}^{int, 2}} \Bigg] \cdot \Delta U_{bl}
    \end{split}
\end{equation}

However, the tensor $\Celastic$ is the material tangent tensor, which is a function of the material model used. This is the information used in the code that each non linear material needs to provide. For that reason, the expression of $\Celastic$ for each non linear material model is described below:

\paragraph{Saint Venant - Kirchhoff:}

\begin{equation*}
    \Celastic = \pderiv{\pkTwo}{\E} = \pderiv{\left( \lambda \, tr(\E) \tensorTwo{\mathrm{I}} + 2\mu \,\E \right)}{\E} \Longrightarrow
\end{equation*}
\begin{equation} \label{eq:C_stk}
    \mathcal{C}_{ijkl} = \lambda \, \delta_{ij}\delta_{kl} + \mu \left( \delta_{ik}\delta_{jl} + \delta_{il}\delta_{jk} \right)
\end{equation}

\paragraph{Compressible Neo-Hookean:}

\begin{equation*}
    \Celastic = \pderiv{\pkTwo}{\E} = \pderiv{\left( \mu \, \left(\tensorTwo{\mathrm{I}} - \C^{-1} \right) - \lambda \, \ln{\J} \, \C^{-1} \right)}{\E} \Longrightarrow
\end{equation*}
\begin{equation} \label{eq:C_nhk}
    \mathcal{C}_{ijkl} = \left( \mu - \lambda \, \ln{\J} \right) \cdot \left( \mathrm{C}_{ik}^{-1} \cdot \mathrm{C}_{jl}^{-1} + \mathrm{C}_{il}^{-1} \cdot \mathrm{C}_{jk}^{-1} \right) + \lambda \cdot \mathrm{C}_{ij}^{-1} \cdot \mathrm{C}_{kl}^{-1}
\end{equation}

\subsection{External virtual work}

{\color{red}\textbf{To do}}


\include{sections/07_code}

\end{document}
